\section{A finite cover for subspaces of controlled height}\label{sec:low-height}

Lemmas~\ref{lem:robust-extraction} and \ref{lem:robust-height} allow us
to pass from a large subset of \(\bits^m\) to a robust subset, with loss
at most \(\rho N\) in its base-two logarithmic size for
\(\rho=\ell^{-D_1}\) and fixed \(D_1>0\), whose difference span \(W\)
satisfies \(H(W)\le \exp(O(N\ell))\). This span may depend on the matrix
under study, so a rank estimate for one fixed space is insufficient.
Theorem~\ref{thm:low-height-cover} replaces all such spans by a small
fixed list while changing the relevant ranks only slightly.
Proposition~\ref{prop:typical-ranks} uses this cover to obtain simultaneous
rank estimates.

For rational subspaces \(U,V\subseteq\R^m\), define their \emph{rank
distance} by
\[
 d(U,V)=\dim U+\dim V-2\dim(U\cap V).
\]
Thus \(d(U,V)\) is the sum of the codimensions of \(U\cap V\) in \(U\)
and \(V\). Recall that \(H(U)=\det(U\cap\Z^m)\), with \(H(0)=1\).
One also has \(d(U^\perp,V^\perp)=d(U,V)\). Passing through the common
subspace \(U^\perp\cap V^\perp\) shows that the nullities of a fixed
symmetric bilinear form restricted to \(U^\perp\times U^\perp\) and
\(V^\perp\times V^\perp\) differ by \(O(d(U,V))\); the same holds for
the ranks of columns projected to these spaces. This is the relation
needed to transfer the fixed-space rank estimates.

\begin{theorem}[Low-height cover]\label{thm:low-height-cover}
Fix \(D,A>0\). For all sufficiently large \(N\), and every integer
\(0\le m\le N\), there is a list \(\mathcal C_{m,N}\) of rational subspaces of
\(\R^m\), depending also on \(D,A\), such that
\[
 |\mathcal C_{m,N}|\le 2^{N^2\ell^{-A}}.
\]
Every rational subspace \(W\subseteq\R^m\) satisfying
\(H(W)\le\exp(N\ell^D)\) has
\[
 \min_{U\in\mathcal C_{m,N}}d(U,W)\le N\ell^{-A}.
\]
\end{theorem}

\begin{proof}
Choose a fixed \(B_0>D+A+2\), and put
\[
 a=\ell^{-B_0},\qquad h(U)=\log_N H(U),\qquad
 \Delta=\frac{N\ell^D}{a\log N}.
\]
Throughout the proof \(N\) is sufficiently large in terms of these
fixed parameters. In particular \(a=o(1)\) and \(1/\log N=o(a)\).
We will encode a space within distance \(\Delta\) of each eligible
\(W\), at a cost \(o(N^2\ell^{-A})\) bits.
We follow the coordinate equations of \(W\) by a nested chain of
minimizing spaces. The height budget limits the total dimension lost in
drops larger than one. A flag in the normal lattice encodes each drop,
and the minimizing inequalities ensure that few flag positions have a
large encoding cost. We then sum these costs over the chain.

\medskip\noindent\textbf{The minimizing chain.}
Let \(r=\dim W\). Choose \(r\) pivot coordinates on which the coordinate
projection of \(W\) is an isomorphism, and write \(W\) as the graph of
rational linear functions on these coordinates. Order the nonpivot
coordinates, and impose their graph equations one at a time. This gives
rational spaces
\[
 \R^m=P_0\supset P_1\supset\cdots\supset P_{m-r}=W.
\]
At a step with coordinate vector \(e\), one has
\(P_{s-1}=P_s\oplus\R e\). Every coordinate vector whose equation has
not yet been imposed belongs to \(P_s\).

For each \(s\), minimize
\[
 F_s(U)=h(U)+a\,d(U,P_s)
\]
over rational subspaces \(U\), write the minimum as \(E_s\), and let
\(U_s\) be the greatest minimizer under inclusion. We justify these
choices. By Lemma~\ref{lem:height-submodularity}, \(h(U)\ge0\), and only
finitely many rational spaces satisfy \(H(U)\le N^{am}\). Testing
\(U=0\) shows that a minimizer can be sought among those spaces, so the
minimum exists.
Lemma~\ref{lem:height-submodularity} and the inequality
\[
 d(U+V,P)+d(U\cap V,P)\le d(U,P)+d(V,P)
\]
show that the sum and intersection of two minimizers for \(F_s\) are
again minimizers. The displayed inequality follows from the dimension
identity for \(U,V\) and the inclusion
\((U\cap P)+(V\cap P)\subseteq (U+V)\cap P\).
Since the set of minimizers is finite, their sum is the greatest one.
Also \(U_0=\R^m\) and \(E_0=0\).

The minimizers descend with the graph spaces. In fact, if \(P'\subseteq
P\), then
\[
 d(U+V,P)+d(U\cap V,P')\le d(U,P)+d(V,P').
\]
Here the required inequality between intersection dimensions follows
from
\((U\cap P)+(V\cap P')\subseteq (U+V)\cap P\), whose two summands
intersect in \(U\cap V\cap P'\). Combining this inequality with
Lemma~\ref{lem:height-submodularity}, for \(U=U_{s-1}\) and \(V=U_s\),
gives
\[
 F_{s-1}(U_{s-1}+U_s)+F_s(U_{s-1}\cap U_s)
 \le E_{s-1}+E_s.
\]
Minimality gives the reverse inequality. Hence \(U_{s-1}+U_s\) is an
old minimizer, and greatestness implies \(U_s\subseteq U_{s-1}\).

Every unprocessed coordinate vector \(e\in P_s\) belongs to \(U_s\).
Indeed, Lemma~\ref{lem:height-submodularity}, together with
\(H(\R e)=1\) and \(H(U\cap\R e)\ge1\), gives
\[
 H(U+\R e)\le H(U).
\]
If \(e\notin U\) but \(e\in P_s\), adjoining \(e\) decreases
\(d(U,P_s)\) by one, contradicting minimality. The energies are
nondecreasing as well. For a step \(P_{s-1}=P_s\oplus\R e\), one has
\[
 d(U+\R e,P_{s-1})\le d(U,P_s).
\]
This follows directly by comparing the intersections after adjoining
\(e\): the new intersection contains
\((U\cap P_s)\oplus\R e\), while the dimensions of \(U\) and \(P_s\)
each increase by at most one. The height inequality just proved,
applied to \(U=U_s\), yields
\(E_{s-1}\le E_s\).

Testing the last objective at \(W\) gives \(E_{m-r}\le h(W)\). Therefore
\[
 d(U_{m-r},W)\le \frac{h(W)}a\le\Delta.
\]
Write \(c_s=\dim U_{s-1}-\dim U_s\) for the dimension drop at step \(s\),
and let \(z\) count the steps with \(c_s=0\). At such a step the old and
new minimizers coincide and contain the coordinate vector \(e\) being
processed. Their intersection with \(P_s\) has dimension one less than
their intersection with \(P_{s-1}\), so \(E_s-E_{s-1}=a\).
Consequently \(z\le h(W)/a\). Write
\(\dim U_{m-r}=r+\theta\). The distance bound gives
\(|\theta|\le h(W)/a\), while summing the drops gives
\(\sum_s c_s=m-r-\theta\). Since there are \(m-r\) steps,
\[
 \sum_{s:c_s\ge2}(c_s-1)=z-\theta\le \frac{2h(W)}a.
\]
It follows that
\begin{equation}\label{eq:multi-drop-budget}
 \sum_{s:c_s\ge2}c_s\le \frac{4h(W)}a\le4\Delta.
\end{equation}

\medskip\noindent\textbf{Normal flags for a transition.}
We next count the choices for a transition \(U_{\rm old}\supseteq
U_{\rm new}\) of drop \(c\ge1\), assuming \(U_{\rm old}\) is known.
Let \(P_{\rm old}\supset P_{\rm new}\) be the corresponding graph
spaces, and set
\[
 \Lambda=(U_{\rm old}\cap\Z^m)^*
\]
in the Euclidean span \(U_{\rm old}\), and let
\(K=\Lambda\cap U_{\rm new}^{\perp}\). This is a primitive rank \(c\)
sublattice of \(\Lambda\), and it determines \(U_{\rm new}\).
For a rational \(V\subseteq U_{\rm old}\), Lemma~\ref{lem:normal-height}
identifies the determinant of its primitive normal lattice as
\[
 \det(\Lambda\cap V^\perp)=\frac{H(V)}{H(U_{\rm old})}.
\]
If \(V\) has codimension \(t\) in \(U_{\rm old}\), the two rank
distances from \(V\) and \(U_{\rm old}\) to any fixed space differ by at
most \(t\). Comparing \(V\) with the old minimizer, and saturating any
sublattice of \(\Lambda\), therefore gives
\[
 h(V)-h(U_{\rm old})
 \ge a\bigl(d(U_{\rm old},P_{\rm old})-d(V,P_{\rm old})\bigr)
 \ge-at
\]
and hence
\[
 \det L\ge N^{-a\rank L}\qquad(L\subseteq\Lambda).
\]
If \(K'\subseteq K\) is primitive of rank \(t\), it is also primitive
in \(\Lambda\), and is the normal lattice of a rational space
\(V'\) with \(U_{\rm new}\subseteq V'\subseteq U_{\rm old}\).
The normal-height identity and comparison with the new minimizer give
\[
 h(U_{\rm new})-h(V')
 \le a\bigl(d(V',P_{\rm new})-d(U_{\rm new},P_{\rm new})\bigr)
 \le a(c-t).
\]
Thus
\begin{equation}\label{eq:cover-normal-balance}
 \frac{\det K}{\det K'}
 =\frac{H(U_{\rm new})}{H(V')}
 \le N^{a(c-t)}.
\end{equation}
These two bounds are the only restrictions on the transition used in
the count.

Choose a saturated flag
\[
 0=K_0\subset K_1\subset\cdots\subset K_c=K
\]
as follows. In the orthogonally projected quotient \(K/K_{i-1}\),
choose a shortest nonzero vector \(v_i\), put \(\mu_i=\|v_i\|\), and
let \(K_i/K_{i-1}\) be the lattice on its line. A shortest nonzero
vector is primitive, so \(K_i/K_{i-1}=\Z v_i\). Each \(K_i\) is
primitive in \(K\), hence also in \(\Lambda\). Put
\(y_i=\log_N\mu_i\). The determinant identity in
Lemma~\ref{lem:lattice-quotient} and Equation~\eqref{eq:cover-normal-balance}
give, for every suffix,
\begin{equation}\label{eq:cover-flag-suffix}
 \sum_{j=i}^c y_j\le a(c-i+1).
\end{equation}
The classical shortest-vector consequence of Minkowski's convex body
theorem in rank \(u\) is
\(\lambda_1(L)\le C\sqrt u\,(\det L)^{1/u}\), with an absolute \(C\).
Apply it to \(K_{i+u-1}/K_{i-1}\). The vector \(v_i\) is shortest even
in the larger lattice \(K/K_{i-1}\), and hence
\begin{equation}\label{eq:cover-flag-range}
 y_i\le \frac1u\sum_{j=i}^{i+u-1}y_j
       +\frac12\log_Nu+O(1/\log N)
 \qquad(1\le u\le c-i+1).
\end{equation}
Taking a full suffix in this inequality shows
\[
 y_i\le \frac12+a+O(1/\log N).
\]
In particular \(y_i\le1\) for large \(N\).

Call position \(i\) expensive if \(y_i>1/2-3a\). A flag of length
\(c\) has \(O(ac)\) expensive positions, and has none when \(c=1\).
For the latter assertion, Equation~\eqref{eq:cover-flag-suffix} gives
\(y_1\le a<1/2-3a\). For the general assertion, suppose an expensive
position exists, and partition the suffix from the first such position
into ranges, each starting at an expensive position and ending just
before the next one, or at the end. For a range of length \(u\),
Equation~\eqref{eq:cover-flag-range} gives
\[
 \sum_{\text{range}}(y_j-a)
 \ge u\left(\frac12-4a-\frac12\log_Nu-O(1/\log N)\right).
\]
For \(u\le N^{2/5}\), this is at least \(u/5\) for large \(N\).
For every \(u\le N\), it is at least \(-C_1au\), since
\(1/\log N=o(a)\). The total over the ranges is nonpositive by
Equation~\eqref{eq:cover-flag-suffix}. The total length of the short
ranges is thus \(O(ac)\). There are at most \(c/N^{2/5}=o(ac)\) long
ranges, and at most the total short length short ranges. This proves
the asserted count.

\medskip\noindent\textbf{Counting flag vectors.}
We count the vector at each flag position in a quotient determined by
the preceding choices. Put
\[
 \Pi=\Lambda/K_{i-1},
\]
again using orthogonal projection. A rank \(t\) sublattice \(\bar L\)
of \(\Pi\) lifts, together with \(K_{i-1}\), to a rank \(t+i-1\)
sublattice of \(\Lambda\). By Lemma~\ref{lem:lattice-quotient}, the
first normal balance bound, and
\(\det K_{i-1}=\prod_{j<i}\mu_j\le N^c\), it follows that
\begin{equation}\label{eq:cover-projected-determinants}
 \det\bar L\ge
 \frac{N^{-a(t+i-1)}}{\det K_{i-1}}
 \ge N^{-at-2c}.
\end{equation}
Apply Lemma~\ref{lem:det-minimizer} to \(\Pi\) with \(T=N^{-2a}\).
Fix a deterministic choice of its minimizing primitive sublattice
\(L_0\), so that this choice requires no additional encoding once
\(\Pi\) is known. If \(j_0=\rank L_0\), the rank-zero competitor and
Equation~\eqref{eq:cover-projected-determinants} imply
\[
 N^{-aj_0-2c}\le\det L_0\le N^{-2aj_0},
 \qquad j_0\le \frac{2c}{a}.
\]
All sublattices of the quotient \(\Omega=\Pi/L_0\) have determinant
at least \(N^{-2a}\) to their rank.

First encode the image of \(v_i\) in \(\Omega\). If \(\Omega\) has
rank zero this image is unique. Otherwise, because its rank is at most
\(N\), Lemma~\ref{lem:gaussian-count} applies at the common smaller
scale
\[
 s_0=\frac{N^{-2a}}{C_2\log(2N)}
\]
for a suitable absolute \(C_2\). For an inexpensive position,
\(\mu_i\le N^{1/2-3a}\), and orthogonal projection does not increase
norm. Lemma~\ref{lem:gaussian-count} therefore bounds the base-two
logarithm of the number of images by
\begin{equation}\label{eq:cover-inexpensive-cost}
 O\bigl(1+N^{1-2a}\log^2N\bigr).
\end{equation}
For an expensive position the general bound on \(y_i\) permits a
common radius \(R=C_3N^{1/2+a}\), with an absolute \(C_3\). Use the
Gaussian comparison in Lemma~\ref{lem:gaussian-count} from \(s_0\)
to \(s_1=R/\sqrt N\). These satisfy \(s_1\ge s_0\) for large \(N\).
Writing \(k=\rank\Omega\le N\), the base-two logarithm of the count is
at most
\[
 O(1)+O(R^2/s_1^2)+k\log_2(s_1/s_0)
 \le O(N)+O\bigl(N(a\log N+\log\log N)\bigr)
 =O(aN\log N).
\]
The last equality uses \(a\log N\gg\log\log N\).

It remains to count the vectors above one fixed image in \(\Omega\).
Use the common radius \(R=C_3N^{1/2+a}\) for both kinds of position,
and let \(\mathcal V\) be the vectors \(v\in\Pi\) above that image with
\(\|v\|\le R\). If \(\mathcal V\) is nonempty, its differences lie in
\(L_0\). Let \(F\) be their real span, of dimension
\[
 j\le j_0\le 2c/a,\qquad j\le N.
\]
The intersection lattice \(L_F=\Pi\cap F\) has rank \(j\), and
Equation~\eqref{eq:cover-projected-determinants} gives
\(\det L_F\ge N^{-aj-2c}\). If \(j=0\), there is at most one vector.
Otherwise choose \(j\) independent differences in \(\mathcal V\),
each of norm at most \(2R\), and let \(L'\subseteq L_F\) be their
integer span. Fix \(v_0\in\mathcal V\). Within each coset of \(L'\)
in \(v_0+L_F\), the translates at the points of \(\mathcal V\) of a
half-open fundamental parallelepiped for these differences are disjoint
up to boundaries. They lie in the \(j\)-dimensional ball centered at
\(v_0\) of radius \(2(1+j)R\). Writing \(B_j(r)\) for the Euclidean
\(j\)-ball of radius \(r\), we may sum over the
\([L_F:L']=\det L'/\det L_F\) cosets to obtain
\[
 |\mathcal V|
 \le \frac{\operatorname{vol}_j B_j(2(1+j)R)}{\det L_F}.
\]
Since \(j\le N\) and \(R=O(N^{1/2+a})\), the logarithm of the
numerator is \(O(j\log N)\). The determinant bound consequently gives
\begin{equation}\label{eq:cover-fiber-cost}
 \log_2|\mathcal V|
 \le O((j+c)\log N)
 =O\left(\frac ca\log N\right).
\end{equation}
This bound is uniform over the fixed image and the preceding flag
choices.

The image and fiber choices specify \(v_i\in\Pi\). They also specify
\(K_i\): for any lift \(\widetilde v_i\in\Lambda\),
\[
 K_i=\Lambda\cap\Span(K_{i-1},\widetilde v_i).
\]
Thus the encoding reconstructs the next quotient and ultimately
\(K\), and hence \(U_{\rm new}\). The counts above can be applied
successively to the prefixes of all encodings of actual minimizing
chains; only such prefixes need satisfy the balance bounds.

\medskip\noindent\textbf{Total encoding cost.}
We finish by summing the costs. Specifying the pivot and coordinate
order, the drops, and the expensive positions costs \(O(N\log N)\)
bits. There are at most \(N\) flag positions altogether, because the
sum of all drops is at most \(m\). Put
\[
 S=\sum_{s:c_s\ge2}c_s\le4\Delta
\]
by Equation~\eqref{eq:multi-drop-budget}. The expensive-position
count gives at most \(O(aS)\) such positions over the whole chain;
there are none in drops of size one. For a drop \(c\), the \(c\)
fiber costs in Equation~\eqref{eq:cover-fiber-cost} total
\(O(c^2\log N/a)\). Using
\(\sum_{s:c_s\ge2}c_s^2\le S^2\), the full base-two logarithmic cost is
therefore bounded as follows. The four terms in the first bound account,
respectively, for the choices of pivots, order, drops, and positions;
for inexpensive images; for expensive images; and for fibers:
\[
 \begin{split}
 O(N\log N)
 &+O\bigl(N(1+N^{1-2a}\log^2N)\bigr)
   +O(a^2N\log N\,S)\\
 &+O\left(\frac{\log N}{a}(S^2+N)\right)\\
 &\le
 O\bigl(N^2N^{-2a}\log^2N\bigr)
 +O(aN^2\ell^D)
 +O\left(\frac{N^2\ell^{2D}}{a^3\log N}\right)
 +O\left(\frac{N\log N}{a}\right).
 \end{split}
\]
Each term is \(o(N^2\ell^{-A})\). For the first,
\(N^{-2a}\log^2N=\exp(-2\log N/\ell^{B_0}+2\ell)
=o(\ell^{-A})\). The second has factor
\(a\ell^D=\ell^{D-B_0}=o(\ell^{-A})\). For the third,
\(\ell^{2D}/(a^3\log N)=\ell^{2D+3B_0}/\log N=o(\ell^{-A})\);
the fourth follows from
\(\ell^{B_0}\log N/N=o(\ell^{-A})\). All exponents here are fixed
before \(N\) tends to infinity.

Collect the terminal spaces decoded by all these encodings, removing
duplicates. This is a list independent of \(W\), of size at most
\(2^{N^2\ell^{-A}}\) for sufficiently large \(N\). Every eligible
\(W\) contributes an encoding whose terminal space satisfies
\[
 d(U_{m-r},W)\le\Delta
 =\frac{N\ell^{D+B_0}}{\log N}
 =o(N\ell^{-A}).
\]
The required distance bound follows, uniformly for \(m\le N\).
\end{proof}
